\documentclass[10pt]{article}
\usepackage[a4paper,margin=26mm]{geometry}
\usepackage{amsmath,amssymb,amsthm,mathrsfs,hyperref}
\newtheorem{theorem}{Theorem}
\newtheorem{lemma}{Lemma}
\newcommand{\PSH}{\operatorname{PSH}}
\newcommand{\FS}{\operatorname{FS}}
\newcommand{\OO}{\mathcal O}
\newcommand{\vol}{\operatorname{vol}}
\title{Complete moving systems, a trace transfer, and relative dual coercivity}
\author{An unreviewed proof bridge for the original big-class equation}
\date{6 October 2026}
\begin{document}
\maketitle
\begin{abstract}
We propose a reconstruction of three missing bridges in the actual big
rational-class equation programme. Complete moving systems are cofinal
in the ordinary positive intersection product, including tests by every
fixed Cartier divisor. Their generic fixed factors are removed on actual
models without changing the original carrier. A filtration trace bound
retains the missing section rank until the auxiliary degree limit is
taken. It transfers a finite-model quadratic contraction estimate to
the genuine relative finite-energy cone. A countable convex-series
argument then proves equivalence of the intrinsic finite-dual condition
and integration of every relative finite-energy test, together with
coercivity. The genuine closure and section-volume energy foundations
are explicit conditional inputs here; the entire equation is not proved.
\end{abstract}

\paragraph{Fixed domain and explicit foundation boundary.}
Let $X$ be every smooth integral projective $n$-fold over an algebraically
closed trivially valued field of characteristic zero. Let $L$ be an actual
big rational bundle, $K=X^{\rm an}$, and $V=\vol_X(L)>0$.
The following genuine foundations, proposed in the relative-floor and
section-volume reconstruction, are \emph{assumptions of this bridge},
to be reconstructed before any original-scope acceptance:

\begin{enumerate}
\item The genuine numerical cone has actual all-point decreasing FS
definitions, convex combinations, decreasing closure, the envelope
property and the normalized reference $V_\theta=P_\theta(0)$.
Every generically normalized genuine $u$ satisfies $u\leq V_\theta$.
\item For every continuous original obstacle $f$, actual section weights
$W_q(f;s)=\inf_K(qf-\log|s|)$ give the raw weight-mean energy
$\mathscr E(f)$. For a rational model obstacle $f_D$, the actual
relative complete-system maxima $\phi_{m,D}$ lie below its genuine
envelope $u_D$ and increase to it at the generic valuation.
The normalized unweighted maxima $\phi_{m,0}$ lie below $V_\theta$.
\item On the genuine cone,
$E(u)=\inf_{f\in C(K),\,f\geq u}\mathscr E(f)$ is concave,
$E(u+c)=E(u)+cV$, $E(V_\theta)=0$, and it is continuous on
decreasing genuine nets. Rational model-majorant envelopes are cofinal
and decrease to each genuine $u$, with their energies decreasing to $E(u)$.
\end{enumerate}
Write $\mathcal E^1=\{u:E(u)>-\infty\}$ and $B(u)=-E(u)$ for
generic-normalized $u$. These assumptions do not include a moment
estimate, mixed products of unbounded potentials, target coercivity,
attainment, comparison or uniqueness.

The ordinary algebraic background is the positive-product construction
of Boucksom--Favre--Jonsson on projective birational models: nef products
are ordered in the pseudoeffective cone, the products below a fixed big
class have a least upper bound and converge in the weak Weil-class
topology, and $\vol(\alpha)=\langle\alpha^d\rangle$ in dimension $d$.
Its purely algebraic characteristic-zero domain is retained. Ordinary
ample global generation, resolution of ideals and flat full-section
asymptotics are also retained. No nonample equation is a background input.

\section{Cofinality, with actual effective remainders}
Let $Z$ be smooth projective of dimension $d$ and $N$ an actual big
rational bundle. Resolve the base ideal of its complete system $|mN|$:
\[
 \mu_m^*(mN)=M_m+G_m,\qquad M_m\text{ globally generated},\quad
 G_m\text{ effective},\qquad A_m=M_m/m.                \tag{1}
\]
The index may be a cofinal factorial divisibility sequence. On common
higher models $\mathfrak a_m^\ell\subset\mathfrak a_{m\ell}$ gives
$A_{m\ell}\geq A_m$ with an actual effective Cartier remainder.
Every $A_m$ is an admissible nef approximation below $\mu_m^*N$.

We prove that these complete systems recover the entire ordinary
positive product, not only its top degree. Take finitely many nef
competitors $\beta_i\leq\mu^*N$ on a common projective model. It
suffices first to take the rational competitors in the defining
positive-product family. Choose an actual rational ample $Q_0$ on
this common model sufficiently small that $\mu^*N-Q_0$ is big.
For small rational $\epsilon>0$,
\[
 Q_i=(1-\epsilon)\beta_i+\epsilon c_1(Q_0)
\]
is ample, and $\mu^*c_1(N)-Q_i$ is big: it is a positive convex
combination of a pseudoeffective class and a big class. Choose an
actual rational bundle representing $Q_i$. Its actual difference
from $\mu^*N$ is big, so after an integral multiple it has a nonzero
section. Thus for one common sufficiently divisible $q$ we have
\[
 q\mu^*N=qQ_i+R_i,\qquad R_i\geq0,
       \qquad qQ_i\text{ globally generated for every }i.          \tag{2}
\]
These are actual bundle identities, including any numerically trivial
differences. Multiplying sections by the actual $R_i$ sections injects
the generated $qQ_i$ systems into the complete $qN$ system.
On a simultaneous resolution its base divisor therefore satisfies
$G_q\leq R_i$, whence $A_q\geq Q_i$ for every $i$.
Nef-product monotonicity gives $A_q^p\geq Q_1\cdots Q_p$.
Letting $\epsilon\downarrow0$ reaches each product of nef competitors.
The least upper bound property and the directed convergence theorem
therefore imply
\[
 A_m^p\longrightarrow\langle N^p\rangle,
                      \qquad 1\leq p\leq d.             \tag{3}
\]
All products are understood as compatible Weil numerical classes.
For $p=d-1$ this includes pairing with every fixed Cartier divisor
on any fixed birational model. A signed test is written as a difference
of two ample Cartier classes; order is paired with the nef tests
separately, never directly with a nonnef central component.
For $p=d$, (3) gives $A_m^d\to\vol(N)$.
The identification for nonrational numerical competitors follows by
the same ample approximation and continuity; it does not change the
actual bundles used in (2).

\section{The actual generic fixed factor and central moving model}
Fix a rational model obstacle $f_D$ for $L$ on a smooth equivariant
compactification $Y\to X\times\mathbf P^1$,
with central fiber $F=\sum b_EE$ and rational vertical Cartier $D$.
Choose one integer $c$ large enough that every relative lattice weight
$w_{mj}+mc$ is nonnegative. This is uniform in $m$: if $-MF\leq D\leq MF$,
the actual section lattice lies between $t^{mM}H^0(X,mL)[t]$ and
$t^{-mM}H^0(X,mL)[t]$, so $c>M$ suffices. Its full absolute basis is
$t^{-(w_{mj}+mc)+r}s_{mj}$, $0\leq r\leq w_{mj}+mc$.
It generates the relative complete lattice at zero and the complete
generic system at infinity. Put $\alpha_D(c)=\pi_X^*L+D+cF$.

Resolve the complete original $|mL|$ on $p_m:X_m\to X$:
\[
 p_m^*L=A_m+T_m,\qquad mA_m\text{ globally generated},\quad T_m\geq0.
                                                               \tag{4}
\]
Pull the total model to a common modification of $X_m\times\mathbf P^1$
and principalize its full global base ideal. Every global section is
divisible by the same horizontal Cartier factor $mT_m$, because its
generic sections belong to the complete original system. After that
division the residual ideal is the unit ideal off time zero: the
relative lattice generates the original complete system there, and
the moving generic $mA_m$ system is free. The residual ideal is
therefore vertical. On its resolution let its effective divisor be
$G_{m,D}^{\rm vert}$. The complete moving bundle divided by $m$ is
\[
 C_m(D)=\rho^*A_m+D+cF-G_{m,D}^{\rm vert}/m.             \tag{5}
\]
It is globally generated and nef. Removing the time twist gives an
actual relatively generated model for $A_m$, with bounded model
potential $\psi_{m,D}$. On $X_m^{\rm an}$ the pullback of the original
FS function is
\[
 \phi_{m,D}\circ p_m^{\rm an}=q_{T_m}+\psi_{m,D},
             \qquad \phi_{m,0}\circ p_m^{\rm an}=q_{T_m}.           \tag{6}
\]
This is an equality after pullback. No identification of supported
semivaluations on the two carriers, and no division at a pole of
$q_{T_m}$, is asserted. For $D=0$ its residual central ideal is trivial.

Apply (3) both to $L$ and to the actual big compactified
$\alpha_D(c)$. It gives
\[
 A_m^n\to V,\qquad C_m(D)^{n+1}\to\vol_Y(\alpha_D(c)),\quad
 C_m(D)^n\to\langle\alpha_D(c)^n\rangle.                \tag{7}
\]
The last limit is against every fixed Cartier test, including the
central components on the fixed model. The finite relative nef model
energy, defined by actual intersections, is
\[
 e_m(\psi_{m,D})=\frac{C_m(D)^{n+1}}{n+1}-cA_m^n
                       \longrightarrow E(u_D).         \tag{8}
\]
Here the section-volume intercept is
$\vol_Y(\alpha_D(c))=(n+1)cV+(n+1)\mathscr E(f_D)$ and
$\mathscr E(f_D)=E(u_D)$ by the stated foundation.
Neither the physical original $\phi_{m,D}$ nor its horizontal factor
is asserted to be in the original $\mathcal E^1$.

\section{A trace inequality with the missing rank retained}
Fix $m$, a bounded relative nef model potential $\psi$ on the actual
$A_m$, and an original continuous obstacle $f$ such that
\[
                    f\circ p_m^{\rm an}\geq q_{T_m}+\psi.          \tag{9}
\]
For divisible auxiliary section degree $\ell$, multiplication by
the actual $\ell T_m$ section injects
$H^0(X_m,\ell A_m)$ into $H^0(X,\ell L)$, by proper birational
normality. Analytification of $p_m$ is surjective, so original
filtration weights can be computed after pullback. For each nonzero
auxiliary section $s$,
\[
 W_\ell(f;s\,s_{\ell T_m})\geq W_\ell(\psi;s).          \tag{10}
\]
Where the product section vanishes, the left bracket is $+\infty$;
no subtraction of two infinities is used. Elsewhere (10) is just (9).
Every section of the original full space also has weight at least
$\ell\min f$, since its canonical section norm is at most one.

Let $r_\ell=h^0(X,\ell L)$ and $d_\ell=h^0(X_m,\ell A_m)$.
Extend an adapted basis of the injected filtered subspace to the
full section space. The trace of its filtration is at least the trace
on that subspace plus $(r_\ell-d_\ell)\ell\min f$.
This follows directly by ordering the finite filtration jumps, or
by the induced filtration on the quotient. Taking normalized traces,
including the vanishing floor error, and then $\ell\to\infty$ for
this \emph{fixed} $m$, yields
\[
 \mathscr E(f)\geq e_m(\psi)+(\min f)(V-A_m^n).           \tag{11}
\]
The auxiliary trace limit equals the relative nef model intersection
energy: its time-twisted bundle is nef and globally generated after
a sufficiently large fixed twist, and the same lattice section count
identifies the linear volume intercept. This also works if the
generic nef bundle has zero top degree.
The order in (11) is auxiliary degree first, fixed original obstacle
next, moving-system index next, and infimum over original majorants
last. No varying $\min f$ error is discarded.

\section{Quadratic contraction from finite positive intersections}
For a nonpositive bounded model potential $\psi$ on a relatively nef
model with generic nef bundle $A$, put $V_A=A^n$, $g=\psi(v_0)\leq0$,
and $B_A(\psi)=-e_A(\psi)$. On a common SNC model define
\[
 a_k=\int(-\psi)\,d\mu(\psi^{[k]},0^{[n-k]})\geq0,
                         \quad 0\leq k\leq n.           \tag{12}
\]
These are finite mixed model intersection measures with their actual
central multiplicities, not products of unbounded genuine potentials.
Positivity follows because all restrictions to the projective central
components are nef. The trivial reference model has measure
$V_A\delta_{v_0}$: exceptional components have zero pulled-back top
intersection. Thus $a_0=-gV_A$.
The finite intersection cocycle and its binomial expansion give
\[
 B_A(r\psi)=\frac r{n+1}\sum_{j=0}^n\sum_{k=0}^j
    \binom jk r^k(1-r)^{j-k}a_k,
             \qquad B_A(\psi)=\frac1{n+1}\sum_{k=0}^n a_k.         \tag{13}
\]
For $k\geq1$, $r^{k+1}\leq r^2$ and
$\sum_{j=k}^n\binom jk=\binom{n+1}{k+1}\leq2^{n+1}$.
The $k=0$ terms are at most $ra_0$. Hence
\[
 B_A(r\psi)\leq r(-g)V_A+2^{n+1}r^2B_A(\psi).          \tag{14}
\]
First take rational $0\leq r\leq1$ so all models are actual rational
bundles. The finite polynomial identity extends it to real $r$.

Let $u=u_D$ be a generically normalized model envelope. Then
$\psi_{m,D}\leq0$ and $g_m=\psi_{m,D}(v_0)\to0$ by the stated
generic section convergence. On the original carrier,
\[
 (1-r)\phi_{m,0}+r\phi_{m,D}
                          \leq(1-r)V_\theta+ru=:u^{(r)}.           \tag{15}
\]
Its pullback is $q_{T_m}+r\psi_{m,D}$, so (11) applies to every
fixed continuous original $f\geq u^{(r)}$. Use (7)--(8) and (14)
while keeping this $f$ fixed, then infimize over all such $f$:
\[
                   B(u^{(r)})\leq C_n r^2B(u),
                            \qquad C_n=2^{n+1}.         \tag{16}
\]
The generic-value error $r(-g_m)A_m^n$ tends to zero. In particular
(16) did not treat a fixed-system physical FS function as energy finite.

For general normalized $u\in\mathcal E^1$, take decreasing model
envelopes $u_i\downarrow u$ and let $c_i=u_i(v_0)\downarrow0$.
Apply (16) to $u_i-c_i$. The unnormalized mixtures
$(1-r)V_\theta+ru_i$ decrease to $u^{(r)}$. Their energy bounds include
the explicit shift $r c_iV$ and costs $-E(u_i)+c_iV$.
Decreasing energy continuity and $c_i\to0$ prove (16) for $u$.
Rational $r$ extends to real interior $r$ by finite concavity of the
energy along this segment; the two endpoints are immediate.

\section{Intrinsic dual energy and actual coercivity}
Let $\nu$ be a positive Radon measure of raw mass $V$ with
$\nu(\{V_\theta=-\infty\})=0$. Relative differences below are defined
only off that null set. For normalized genuine $u$ put
\[
 I_\nu(u)=\int(V_\theta-u)\,d\nu\geq0,
       \qquad D_\nu=\sup_{u\in\mathcal E^1,\,u(v_0)=0}
                       \{I_\nu(u)-B(u)\}.              \tag{17}
\]
We prove that $D_\nu<\infty$ is equivalent to $I_\nu(u)<\infty$
for every relative finite-energy test. One implication is immediate.
For the other, suppose every such test is integrable. There is then
a finite constant $C_\nu$ with
\[
             I_\nu(u)\leq C_\nu\quad\text{if }B(u)\leq1.          \tag{18}
\]
If not, choose normalized $u_j$ with $B(u_j)\leq1$ and
$I_\nu(u_j)\geq4^j$. The finite convex mixtures
\[
 v_N=\Bigl(1-\sum_{j=1}^N2^{-j}\Bigr)V_\theta
                         +\sum_{j=1}^N2^{-j}u_j
\]
are genuine, decrease, and have generic value zero. Concavity bounds
their costs by one. Decreasing closure and energy continuity give
a normalized $v\in\mathcal E^1$ with $B(v)\leq1$.
On the complement of the null reference pole set,
$V_\theta-v=\sum_j2^{-j}(V_\theta-u_j)$.
Monotone integration gives $I_\nu(v)=\infty$, a contradiction.

For $B(u)\geq1/C_n$ choose $r=(C_nB(u))^{-1/2}$ in (16).
Then $B(u^{(r)})\leq1$ and
\[
                 I_\nu(u)\leq C_\nu\sqrt{C_nB(u)}.      \tag{19}
\]
For smaller costs (18) applies. Therefore the supremum in (17) is
finite. Conversely, if $D_\nu$ is finite, apply its defining bound
to $u^{(r)}$ and use $I_\nu(u^{(r)})=rI_\nu(u)$ and (16):
\[
 I_\nu(u)\leq C_n rB(u)+D_\nu/r.
\]
For every $0<\epsilon\leq1$ take $r=\epsilon/C_n$ to obtain
\[
 I_\nu(u)\leq\epsilon B(u)+C_nD_\nu/\epsilon,
 \qquad E(u)+I_\nu(u)
       \leq-(1-\epsilon)B(u)+C_nD_\nu/\epsilon.           \tag{20}
\]
With $\epsilon<1$ this is an actual coercivity estimate on the entire
normalized genuine cone, rather than a formal duality assertion.

\paragraph{Exact result and remaining work.}
This candidate supplies the complete moving-series positive-product
route, the missing-rank trace bound and, conditional on the explicit
genuine foundations, quadratic contraction and intrinsic target
coercivity. Every new application requires independent reconstruction.
It does not prove that every output measure has null reference poles,
the negligible-discrepancy bridge, the unbounded mixed moment bound,
uniform absolute continuous-test approximation, weighted identification,
strong fixed-target integration, actual all-target attainment, strict
comparison or all-solution uniqueness. Those are still required for
the original equation. No original problem is closed by this draft.

\begin{thebibliography}{9}
\bibitem{BFJ} S. Boucksom, C. Favre and M. Jonsson,
\emph{Differentiability of volumes of divisors and a problem of Teissier},
\href{https://arxiv.org/html/math/0608260v2}{arXiv:math/0608260v2},
Definition 2.5, Proposition 2.3, Lemmas 2.6--2.7, Proposition 2.13
and Theorem 3.1. The stated algebraic characteristic-zero domain is
retained. Fujita and Hodge-index foundations remain ordinary background.
\bibitem{DEL} J.-P. Demailly, L. Ein and R. Lazarsfeld,
\href{https://arxiv.org/html/math/0002035v2}{A subadditivity property
of multiplier ideals}, Proposition 3.6 and its complete proof:
complete moving top intersections. Our cofinal product proof above
also treats fixed signed Cartier tests.
\bibitem{foundations} Relative-floor closure candidate, 6 October 2026,
and the completed original P16.5 partial section-volume reconstruction.
Their genuine closure/energy statements are explicit conditional
inputs here. The separate compactification and finite-field descent
candidate is under independent review. Historical verdicts and file
hashes supply no mathematical extension of those inputs.
\end{thebibliography}
\end{document}
